% Options for packages loaded elsewhere
% Options for packages loaded elsewhere
\PassOptionsToPackage{unicode}{hyperref}
\PassOptionsToPackage{hyphens}{url}
\PassOptionsToPackage{dvipsnames,svgnames,x11names}{xcolor}
%
\documentclass[
  letterpaper,
  DIV=11,
  numbers=noendperiod]{scrartcl}
\usepackage{xcolor}
\usepackage{amsmath,amssymb}
\setcounter{secnumdepth}{5}
\usepackage{iftex}
\ifPDFTeX
  \usepackage[T1]{fontenc}
  \usepackage[utf8]{inputenc}
  \usepackage{textcomp} % provide euro and other symbols
\else % if luatex or xetex
  \usepackage{unicode-math} % this also loads fontspec
  \defaultfontfeatures{Scale=MatchLowercase}
  \defaultfontfeatures[\rmfamily]{Ligatures=TeX,Scale=1}
\fi
\usepackage{lmodern}
\ifPDFTeX\else
  % xetex/luatex font selection
\fi
% Use upquote if available, for straight quotes in verbatim environments
\IfFileExists{upquote.sty}{\usepackage{upquote}}{}
\IfFileExists{microtype.sty}{% use microtype if available
  \usepackage[]{microtype}
  \UseMicrotypeSet[protrusion]{basicmath} % disable protrusion for tt fonts
}{}
\makeatletter
\@ifundefined{KOMAClassName}{% if non-KOMA class
  \IfFileExists{parskip.sty}{%
    \usepackage{parskip}
  }{% else
    \setlength{\parindent}{0pt}
    \setlength{\parskip}{6pt plus 2pt minus 1pt}}
}{% if KOMA class
  \KOMAoptions{parskip=half}}
\makeatother
% Make \paragraph and \subparagraph free-standing
\makeatletter
\ifx\paragraph\undefined\else
  \let\oldparagraph\paragraph
  \renewcommand{\paragraph}{
    \@ifstar
      \xxxParagraphStar
      \xxxParagraphNoStar
  }
  \newcommand{\xxxParagraphStar}[1]{\oldparagraph*{#1}\mbox{}}
  \newcommand{\xxxParagraphNoStar}[1]{\oldparagraph{#1}\mbox{}}
\fi
\ifx\subparagraph\undefined\else
  \let\oldsubparagraph\subparagraph
  \renewcommand{\subparagraph}{
    \@ifstar
      \xxxSubParagraphStar
      \xxxSubParagraphNoStar
  }
  \newcommand{\xxxSubParagraphStar}[1]{\oldsubparagraph*{#1}\mbox{}}
  \newcommand{\xxxSubParagraphNoStar}[1]{\oldsubparagraph{#1}\mbox{}}
\fi
\makeatother

\usepackage{color}
\usepackage{fancyvrb}
\newcommand{\VerbBar}{|}
\newcommand{\VERB}{\Verb[commandchars=\\\{\}]}
\DefineVerbatimEnvironment{Highlighting}{Verbatim}{commandchars=\\\{\}}
% Add ',fontsize=\small' for more characters per line
\usepackage{framed}
\definecolor{shadecolor}{RGB}{241,243,245}
\newenvironment{Shaded}{\begin{snugshade}}{\end{snugshade}}
\newcommand{\AlertTok}[1]{\textcolor[rgb]{0.68,0.00,0.00}{#1}}
\newcommand{\AnnotationTok}[1]{\textcolor[rgb]{0.37,0.37,0.37}{#1}}
\newcommand{\AttributeTok}[1]{\textcolor[rgb]{0.40,0.45,0.13}{#1}}
\newcommand{\BaseNTok}[1]{\textcolor[rgb]{0.68,0.00,0.00}{#1}}
\newcommand{\BuiltInTok}[1]{\textcolor[rgb]{0.00,0.23,0.31}{#1}}
\newcommand{\CharTok}[1]{\textcolor[rgb]{0.13,0.47,0.30}{#1}}
\newcommand{\CommentTok}[1]{\textcolor[rgb]{0.37,0.37,0.37}{#1}}
\newcommand{\CommentVarTok}[1]{\textcolor[rgb]{0.37,0.37,0.37}{\textit{#1}}}
\newcommand{\ConstantTok}[1]{\textcolor[rgb]{0.56,0.35,0.01}{#1}}
\newcommand{\ControlFlowTok}[1]{\textcolor[rgb]{0.00,0.23,0.31}{\textbf{#1}}}
\newcommand{\DataTypeTok}[1]{\textcolor[rgb]{0.68,0.00,0.00}{#1}}
\newcommand{\DecValTok}[1]{\textcolor[rgb]{0.68,0.00,0.00}{#1}}
\newcommand{\DocumentationTok}[1]{\textcolor[rgb]{0.37,0.37,0.37}{\textit{#1}}}
\newcommand{\ErrorTok}[1]{\textcolor[rgb]{0.68,0.00,0.00}{#1}}
\newcommand{\ExtensionTok}[1]{\textcolor[rgb]{0.00,0.23,0.31}{#1}}
\newcommand{\FloatTok}[1]{\textcolor[rgb]{0.68,0.00,0.00}{#1}}
\newcommand{\FunctionTok}[1]{\textcolor[rgb]{0.28,0.35,0.67}{#1}}
\newcommand{\ImportTok}[1]{\textcolor[rgb]{0.00,0.46,0.62}{#1}}
\newcommand{\InformationTok}[1]{\textcolor[rgb]{0.37,0.37,0.37}{#1}}
\newcommand{\KeywordTok}[1]{\textcolor[rgb]{0.00,0.23,0.31}{\textbf{#1}}}
\newcommand{\NormalTok}[1]{\textcolor[rgb]{0.00,0.23,0.31}{#1}}
\newcommand{\OperatorTok}[1]{\textcolor[rgb]{0.37,0.37,0.37}{#1}}
\newcommand{\OtherTok}[1]{\textcolor[rgb]{0.00,0.23,0.31}{#1}}
\newcommand{\PreprocessorTok}[1]{\textcolor[rgb]{0.68,0.00,0.00}{#1}}
\newcommand{\RegionMarkerTok}[1]{\textcolor[rgb]{0.00,0.23,0.31}{#1}}
\newcommand{\SpecialCharTok}[1]{\textcolor[rgb]{0.37,0.37,0.37}{#1}}
\newcommand{\SpecialStringTok}[1]{\textcolor[rgb]{0.13,0.47,0.30}{#1}}
\newcommand{\StringTok}[1]{\textcolor[rgb]{0.13,0.47,0.30}{#1}}
\newcommand{\VariableTok}[1]{\textcolor[rgb]{0.07,0.07,0.07}{#1}}
\newcommand{\VerbatimStringTok}[1]{\textcolor[rgb]{0.13,0.47,0.30}{#1}}
\newcommand{\WarningTok}[1]{\textcolor[rgb]{0.37,0.37,0.37}{\textit{#1}}}

\usepackage{longtable,booktabs,array}
\newcounter{none} % for unnumbered tables
\usepackage{calc} % for calculating minipage widths
% Correct order of tables after \paragraph or \subparagraph
\usepackage{etoolbox}
\makeatletter
\patchcmd\longtable{\par}{\if@noskipsec\mbox{}\fi\par}{}{}
\makeatother
% Allow footnotes in longtable head/foot
\IfFileExists{footnotehyper.sty}{\usepackage{footnotehyper}}{\usepackage{footnote}}
\makesavenoteenv{longtable}
\usepackage{graphicx}
\makeatletter
\newsavebox\pandoc@box
\newcommand*\pandocbounded[1]{% scales image to fit in text height/width
  \sbox\pandoc@box{#1}%
  \Gscale@div\@tempa{\textheight}{\dimexpr\ht\pandoc@box+\dp\pandoc@box\relax}%
  \Gscale@div\@tempb{\linewidth}{\wd\pandoc@box}%
  \ifdim\@tempb\p@<\@tempa\p@\let\@tempa\@tempb\fi% select the smaller of both
  \ifdim\@tempa\p@<\p@\scalebox{\@tempa}{\usebox\pandoc@box}%
  \else\usebox{\pandoc@box}%
  \fi%
}
% Set default figure placement to htbp
\def\fps@figure{htbp}
\makeatother





\setlength{\emergencystretch}{3em} % prevent overfull lines

\providecommand{\tightlist}{%
  \setlength{\itemsep}{0pt}\setlength{\parskip}{0pt}}



 


\KOMAoption{captions}{tableheading}
\makeatletter
\@ifpackageloaded{tcolorbox}{}{\usepackage[skins,breakable]{tcolorbox}}
\@ifpackageloaded{fontawesome5}{}{\usepackage{fontawesome5}}
\definecolor{quarto-callout-color}{HTML}{909090}
\definecolor{quarto-callout-note-color}{HTML}{0758E5}
\definecolor{quarto-callout-important-color}{HTML}{CC1914}
\definecolor{quarto-callout-warning-color}{HTML}{EB9113}
\definecolor{quarto-callout-tip-color}{HTML}{00A047}
\definecolor{quarto-callout-caution-color}{HTML}{FC5300}
\definecolor{quarto-callout-color-frame}{HTML}{acacac}
\definecolor{quarto-callout-note-color-frame}{HTML}{4582ec}
\definecolor{quarto-callout-important-color-frame}{HTML}{d9534f}
\definecolor{quarto-callout-warning-color-frame}{HTML}{f0ad4e}
\definecolor{quarto-callout-tip-color-frame}{HTML}{02b875}
\definecolor{quarto-callout-caution-color-frame}{HTML}{fd7e14}
\makeatother
\makeatletter
\@ifpackageloaded{caption}{}{\usepackage{caption}}
\AtBeginDocument{%
\ifdefined\contentsname
  \renewcommand*\contentsname{Table of contents}
\else
  \newcommand\contentsname{Table of contents}
\fi
\ifdefined\listfigurename
  \renewcommand*\listfigurename{List of Figures}
\else
  \newcommand\listfigurename{List of Figures}
\fi
\ifdefined\listtablename
  \renewcommand*\listtablename{List of Tables}
\else
  \newcommand\listtablename{List of Tables}
\fi
\ifdefined\figurename
  \renewcommand*\figurename{Figure}
\else
  \newcommand\figurename{Figure}
\fi
\ifdefined\tablename
  \renewcommand*\tablename{Table}
\else
  \newcommand\tablename{Table}
\fi
}
\@ifpackageloaded{float}{}{\usepackage{float}}
\floatstyle{ruled}
\@ifundefined{c@chapter}{\newfloat{codelisting}{h}{lop}}{\newfloat{codelisting}{h}{lop}[chapter]}
\floatname{codelisting}{Listing}
\newcommand*\listoflistings{\listof{codelisting}{List of Listings}}
\makeatother
\makeatletter
\makeatother
\makeatletter
\@ifpackageloaded{caption}{}{\usepackage{caption}}
\@ifpackageloaded{subcaption}{}{\usepackage{subcaption}}
\makeatother
\usepackage{bookmark}
\IfFileExists{xurl.sty}{\usepackage{xurl}}{} % add URL line breaks if available
\urlstyle{same}
\hypersetup{
  pdftitle={Direct DC Integration of Wind Turbines and PEM Electrolysers},
  colorlinks=true,
  linkcolor={blue},
  filecolor={Maroon},
  citecolor={Blue},
  urlcolor={Blue},
  pdfcreator={LaTeX via pandoc}}


\title{Direct DC Integration of Wind Turbines and PEM Electrolysers}
\usepackage{etoolbox}
\makeatletter
\providecommand{\subtitle}[1]{% add subtitle to \maketitle
  \apptocmd{\@title}{\par {\large #1 \par}}{}{}
}
\makeatother
\subtitle{Why the model assumes direct turbine-DC-link coupling is
technically possible}
\author{}
\date{}
\begin{document}
\maketitle

\renewcommand*\contentsname{Table of contents}
{
\setcounter{tocdepth}{3}
\tableofcontents
}

\section{Purpose of this section}\label{purpose-of-this-section}

The model assumes that, in a decentralised hydrogen architecture, the
PEM electrolyser can be integrated directly on the wind turbine DC link.
In other words, electrical power from the turbine is not exported as
grid-quality AC and then rectified again for electrolysis. Instead, the
turbine-side converter produces a controlled DC link, and the PEM stack
or stack array is connected to that DC system.

The purpose of this section is to explain why that assumption is
reasonable. The argument is made in two steps:

\begin{enumerate}
\def\labelenumi{\arabic{enumi}.}
\tightlist
\item
  direct DC integration has a clear system-level benefit; and
\item
  direct DC integration without a separate electrolyser-side DC/DC
  converter is technically plausible if the turbine converter, PEM
  stack, and balance of plant are designed as one system.
\end{enumerate}

The assumed electrical chain is:

\begin{Shaded}
\begin{Highlighting}[]
\NormalTok{wind generator {-}\textgreater{} active front{-}end rectifier / turbine{-}side converter {-}\textgreater{} controlled DC link {-}\textgreater{} PEM stack or PEM stack array}
\end{Highlighting}
\end{Shaded}

The turbine-side rectifier is assumed to be an \textbf{active front end
(AFE)}. It controls generator torque and DC-link voltage. The PEM stack
or PEM stack array is connected directly to this controlled DC bus.

The architecture considered here has:

\begin{itemize}
\tightlist
\item
  no AC grid export from the turbine;
\item
  no grid-side inverter;
\item
  no separate electrolyser-side DC/DC converter;
\item
  no dedicated stack current regulator between the DC link and the PEM
  stack;
\item
  only protection and auxiliary devices such as DC-link capacitance,
  contactors, precharge circuits, dump load, brake chopper, crowbar, and
  emergency shutdown hardware.
\end{itemize}

\begin{tcolorbox}[enhanced jigsaw, arc=.35mm, bottomrule=.15mm, bottomtitle=1mm, breakable, colback=white, colbacktitle=quarto-callout-note-color!10!white, colframe=quarto-callout-note-color-frame, coltitle=black, left=2mm, leftrule=.75mm, opacityback=0, opacitybacktitle=0.6, rightrule=.15mm, title=\textcolor{quarto-callout-note-color}{\faInfo}\hspace{0.5em}{Model assumption}, titlerule=0mm, toprule=.15mm, toptitle=1mm]

The model assumes that direct DC integration is possible at system
level. This does not imply that a PEM stack can simply be connected to
an uncontrolled DC source. It assumes an integrated design in which the
turbine AFE, DC-link control, PEM stack configuration, module switching,
and balance-of-plant controls jointly keep the stack within its
allowable operating envelope.

\end{tcolorbox}

\section{Why direct DC integration is
beneficial}\label{why-direct-dc-integration-is-beneficial}

The main benefit of direct DC integration is that it removes electrical
conversion stages that are otherwise required only because the energy is
temporarily handled as grid-quality AC.

In a conventional AC-coupled architecture, the electrical chain is
typically:

\begin{Shaded}
\begin{Highlighting}[]
\NormalTok{wind generator {-}\textgreater{} machine{-}side converter {-}\textgreater{} DC link {-}\textgreater{} grid{-}side inverter {-}\textgreater{} transformer {-}\textgreater{} AC collection/export {-}\textgreater{} transformer/rectifier {-}\textgreater{} PEM stack}
\end{Highlighting}
\end{Shaded}

This architecture is logical when the product is electricity, or when
many turbines feed a centralised electrolyser through an AC collection
grid. However, if the only product is hydrogen and there is no AC
export, several parts of the chain become avoidable from a system
perspective.

A PEM electrolyser is fundamentally a DC load. A modern full-converter
wind turbine also already contains a DC link internally. Direct
integration therefore tries to connect two parts of the system that are
already naturally DC:

\begin{Shaded}
\begin{Highlighting}[]
\NormalTok{wind turbine internal DC link {-}\textgreater{} PEM stack DC load}
\end{Highlighting}
\end{Shaded}

This creates four main benefits.

\subsection{Fewer conversion stages}\label{fewer-conversion-stages}

Each electrical conversion stage adds cost, losses, cooling demand,
protection equipment, and controls. Direct DC integration removes the
need to convert the turbine DC link back to grid-quality AC and then
convert that AC back to DC for the electrolyser.

ITM Power's wind-turbine electrolyser concept makes this argument
explicitly. It describes a system with:

\begin{quote}
``an electrolyser placed at the wind turbine \ldots{} directly
electrically connected to the DC link in the wind turbine''
\end{quote}

and states that Type IV turbine generators and their DC links:

\begin{quote}
``have the potential to power the electrolyser directly''
\end{quote}

ITM describes the benefit as:

\begin{quote}
``fewer power conversion steps''
\end{quote}

Source: ITM Power, \emph{Wind turbine electrolyser integration}:
\url{https://itm-power.com/news/wind-turbine-electrolyser-integraion}

This directly supports the modelling idea: if electricity is not the
final export product, it is inefficient to force the system through the
full AC-export chain before feeding a DC electrolyser.

\subsection{Lower electrical infrastructure
requirement}\label{lower-electrical-infrastructure-requirement}

If hydrogen is produced at the turbine, the infield system can shift
from electrical collection to hydrogen collection. This can reduce or
avoid equipment such as turbine grid-side inverters, step-up
transformers, offshore AC collection strings, offshore electrical
platforms, and centralised electrolyser rectifier infrastructure.

The benefit is not only efficiency. It is also structural: the offshore
wind farm no longer needs to be designed primarily as an electrical
export asset. It can be designed as a distributed hydrogen-production
asset.

\subsection{Better fit with a decentralised hydrogen
architecture}\label{better-fit-with-a-decentralised-hydrogen-architecture}

Direct DC integration aligns the electrical design with the value-chain
design. In a decentralised architecture, each turbine becomes a
hydrogen-production unit. The infield network collects hydrogen rather
than electricity.

This avoids moving large amounts of electrical power through an infield
AC grid only to convert it back to DC at a central electrolyser. It also
reduces the relevance of AC string capacity limits and offshore
electrical platform sizing in the model.

\subsection{Potentially lower cost, mass, and
losses}\label{potentially-lower-cost-mass-and-losses}

Removing conversion stages can reduce:

\begin{itemize}
\tightlist
\item
  converter equipment;
\item
  transformers;
\item
  filters;
\item
  switchgear;
\item
  cooling systems;
\item
  electrical platform weight;
\item
  conversion losses;
\item
  interface complexity between turbine and central electrolyser.
\end{itemize}

These benefits are the reason direct DC integration is worth modelling,
even if it requires a more integrated turbine-electrolyser control
system.

\begin{tcolorbox}[enhanced jigsaw, arc=.35mm, bottomrule=.15mm, bottomtitle=1mm, breakable, colback=white, colbacktitle=quarto-callout-tip-color!10!white, colframe=quarto-callout-tip-color-frame, coltitle=black, left=2mm, leftrule=.75mm, opacityback=0, opacitybacktitle=0.6, rightrule=.15mm, title=\textcolor{quarto-callout-tip-color}{\faLightbulb}\hspace{0.5em}{Modelling implication}, titlerule=0mm, toprule=.15mm, toptitle=1mm]

The modelling assumption is not that power electronics disappear. The
assumption is that the turbine-side AFE remains, while the downstream
grid-side inverter and electrolyser-side rectification/DC conversion
stages are not modelled as separate conversion steps for the
decentralised case.

\end{tcolorbox}

\section{Why direct DC integration should be technically
possible}\label{why-direct-dc-integration-should-be-technically-possible}

Direct DC integration should be technically possible because the two key
electrical subsystems are compatible at a basic level:

\begin{itemize}
\tightlist
\item
  a Type IV wind turbine already contains a controlled DC link; and
\item
  a PEM electrolyser stack is a DC electrochemical load.
\end{itemize}

The question is therefore not whether PEM electrolysis can use DC power.
It can. The real question is whether the PEM stack can be operated
safely and efficiently from the turbine DC link \textbf{without adding a
separate DC/DC converter between the DC link and the stack}.

The answer is: yes, if the system is designed around voltage matching
and active DC-link control.

\section{Electrical operating
principle}\label{electrical-operating-principle}

A PEM electrolyser stack draws current according to its polarization
curve. The stack voltage can be represented conceptually as:

\begin{Shaded}
\begin{Highlighting}[]
\NormalTok{V\_stack = N\_cells * V\_cell(j, temperature, pressure, hydration, ageing)}
\end{Highlighting}
\end{Shaded}

where \texttt{j} is current density.

With a dedicated DC/DC converter, the converter normally regulates stack
current. Without that converter, stack current is instead determined by
the controlled DC-link voltage and the PEM polarization curve.

The direct-coupled design therefore becomes a voltage-matched system:

\begin{Shaded}
\begin{Highlighting}[]
\NormalTok{turbine DC{-}link operating range \textasciitilde{}= PEM stack operating voltage range}
\end{Highlighting}
\end{Shaded}

If the PEM stack cell count and module configuration are selected so
that the stack voltage range matches the turbine DC-link operating
range, the stack can be supplied directly from the DC link. The AFE then
acts as the main electrical actuator by regulating the DC-link voltage
and controlling generator torque.

This is why the assumption is credible: the turbine-side AFE can provide
controlled DC power, and the PEM stack can accept controlled DC power.

\section{What replaces the missing DC/DC
converter}\label{what-replaces-the-missing-dcdc-converter}

Removing the electrolyser-side DC/DC converter does not remove the need
for control. It relocates the control function.

In a conventional system, the DC/DC converter or rectifier supply
normally provides:

\begin{itemize}
\tightlist
\item
  stack current control;
\item
  voltage matching;
\item
  current limiting;
\item
  current ripple reduction;
\item
  controlled start-up and shutdown;
\item
  some degree of electrical decoupling;
\item
  fault management and protection.
\end{itemize}

In the direct DC architecture, these functions must be provided by the
integrated system:

\begin{itemize}
\tightlist
\item
  the AFE regulates DC-link voltage;
\item
  generator torque control limits the mechanical-to-electrical power
  conversion;
\item
  pitch control curtails aerodynamic power when required;
\item
  stack modules can be switched in and out to adjust the effective PEM
  load;
\item
  optional series/parallel reconfiguration can adjust stack voltage
  range;
\item
  a DC-link capacitor handles short electrical transients;
\item
  a brake chopper or dump load handles stack rejection and emergency
  overvoltage;
\item
  balance-of-plant interlocks ensure water, temperature, pressure, and
  gas-purity limits are respected.
\end{itemize}

The most credible architecture is therefore:

\begin{Shaded}
\begin{Highlighting}[]
\NormalTok{PMSG wind turbine {-}\textgreater{} AFE {-}\textgreater{} controlled DC link {-}\textgreater{} segmented PEM stack array}
\end{Highlighting}
\end{Shaded}

with module switching and emergency dump/brake protection.

\section{Evidence supporting the
assumption}\label{evidence-supporting-the-assumption}

\subsection{Industrial concept: ITM Power and
Ørsted}\label{industrial-concept-itm-power-and-uxf8rsted}

ITM Power explicitly describes the direct DC-link integration concept.
The proposed system places the electrolyser at the wind turbine and
connects it directly to the turbine DC link:

\begin{quote}
``an electrolyser placed at the wind turbine \ldots{} directly
electrically connected to the DC link in the wind turbine''
\end{quote}

ITM also states that Type IV turbine generators and their DC links:

\begin{quote}
``have the potential to power the electrolyser directly''
\end{quote}

and that the approach enables:

\begin{quote}
``fewer power conversion steps''
\end{quote}

Source: ITM Power, \emph{Wind turbine electrolyser integration}:
\url{https://itm-power.com/news/wind-turbine-electrolyser-integraion}

This source is important because it confirms that direct turbine-DC-link
coupling is not merely a theoretical modelling shortcut. It is an
identified industrial concept for reducing conversion stages in offshore
wind-to-hydrogen systems.

However, ITM also states:

\begin{quote}
``Such a system has not been demonstrated in practice''
\end{quote}

Therefore, this reference supports the feasibility and relevance of the
model assumption, but it does not prove a commercial MW-scale no-DC/DC
implementation.

\subsection{Academic control concept: Jing et
al.}\label{academic-control-concept-jing-et-al.}

Jing et al.~describe a non-grid-connected wind-hydrogen architecture
that:

\begin{quote}
``directly connects the Source wind turbines and the Load hydrogen
production devices through a DC bus''
\end{quote}

The electrolyser load is described as:

\begin{quote}
``a modular proton exchange membrane (PEM) water electrolytic hydrogen
production array''
\end{quote}

The control strategy regulates:

\begin{quote}
``the size of the PEM array \ldots{} and DC voltage collaboratively''
\end{quote}

Source: Jing et al., \emph{Research on the Structure and Control
Strategy of Energy Router for Non-grid-connected Source-Load Hydrogen
Production System}:
\url{https://link.springer.com/content/pdf/10.1007/s42835-021-00774-w.pdf}

This is directly aligned with the model assumption. Instead of using a
separate DC/DC converter to regulate the PEM stack, the system controls
the DC bus and changes the effective PEM array size. This supports the
idea that a segmented stack array can replace part of the function of an
electrolyser-side converter.

The limitation is that this is an academic control and simulation
reference, not evidence of a commercial field system.

\subsection{Small-scale experimental evidence: RMIT / Air
403}\label{small-scale-experimental-evidence-rmit-air-403}

The RMIT direct-coupling work explicitly investigated wind-to-PEM
coupling:

\begin{quote}
``without the need for intervening voltage converters and maximum power
point tracking electronics''
\end{quote}

Source: RMIT, \emph{Direct coupling of a fixed-pitch variable speed
aerogenerator to PEM electrolysers in a wind-hydrogen system for remote
area power supply}:
\url{https://research-repository.rmit.edu.au/articles/conference_contribution/Direct_coupling_of_a_fixed-pitch_variable_speed_aerogenerator_to_PEM_electrolysers_in_a_wind-hydrogen_system_for_remote_area_power_supply/27369558}

This supports the physical feasibility of direct wind-to-PEM coupling
without voltage conversion. It is not directly scalable evidence for a
modern MW offshore turbine, but it demonstrates that matching generator
and PEM stack electrical characteristics can produce a working
direct-coupled system.

\subsection{Counterpoint: mainstream DC wind concepts often retain
DC/DC}\label{counterpoint-mainstream-dc-wind-concepts-often-retain-dcdc}

A recent paper on DC wind-turbine concepts notes that DC coupling is
attractive because:

\begin{quote}
``wind turbine converters and electrolyzers operate internally with dc
voltage''
\end{quote}

and that DC concepts can reduce equipment because:

\begin{quote}
``the number of conversion stages and filters is reduced''
\end{quote}

However, the same paper states:

\begin{quote}
``A key component for the application of dc technology in WTs is a
high-power dc-dc converter''
\end{quote}

with functions including:

\begin{quote}
``stepping up/down dc voltages, galvanic isolation and fault
ride-through capability''
\end{quote}

Source: \emph{Evaluation of wind turbine concepts for coupling wind
energy and electrolysis}:
\url{https://link.springer.com/article/10.1007/s10010-025-00796-0}

This is an important qualification. It shows that direct DC integration
is attractive, but a separate DC/DC converter is still often used in
mainstream DC architectures because it provides voltage adaptation,
isolation, and protection. Therefore, the no-DC/DC assumption is
technically defensible but should be treated as an integrated-design
assumption rather than a standard off-the-shelf configuration.

\section{Technical conditions for
feasibility}\label{technical-conditions-for-feasibility}

\subsection{Voltage matching}\label{voltage-matching}

The PEM stack voltage range must match the turbine DC-link voltage
range. This can be done by selecting the number of cells in series and
by arranging stack modules in suitable series/parallel strings.

If the DC-link voltage is too low, the stack will produce little or no
hydrogen. If it is too high, current density can exceed stack limits.
The design therefore needs voltage matching at rated power and
acceptable operation across the expected wind-power range.

\subsection{Current-density control}\label{current-density-control}

The PEM stack must remain within its allowable current-density range. A
review of dynamic PEM operation notes:

\begin{quote}
``Dynamic operation of PEM electrolyzers via modulation of the operating
current density brings with it several challenges - namely, heat
management, gas crossover and associated safety concerns, and stack
degradation.''
\end{quote}

Source: \url{https://arxiv.org/html/2405.06766v1}

In the direct DC architecture, current density is constrained by DC-link
voltage control, wind curtailment, and stack-module switching rather
than by a dedicated DC/DC converter.

\subsection{Low-load and gas-crossover
limits}\label{low-load-and-gas-crossover-limits}

Low-load operation can create gas-purity and crossover concerns. A
direct wind-coupled system will naturally encounter low-load operation
unless it is curtailed, idled, or shut down below a defined threshold.

A 2025 RSC paper on PEM operation at differential pressure states that
understanding hydrogen permeation is:

\begin{quote}
``critical towards developing effective gas recombination strategies
that enable safe operation''
\end{quote}

Source:
\url{https://pubs.rsc.org/en/Content/ArticleLanding/2025/EE/D5EE00048C}

The model assumption should therefore include a minimum safe operating
region or shutdown logic. Direct DC connection should not be interpreted
as uncontrolled low-current operation at arbitrary pressure.

\subsection{Dynamic response}\label{dynamic-response}

PEM electrolysers are suitable for dynamic renewable operation, but
dynamic operation still affects degradation, thermal control, pressure
control, and gas management. A review of PEM dynamic operation states
that renewable intermittency:

\begin{quote}
``poses several technological challenges and financial questions''
\end{quote}

Source: \url{https://www.mdpi.com/1996-1073/18/9/2154}

This does not invalidate direct DC integration. It means the stack and
balance of plant must be qualified for turbine-driven operation,
including ramps, starts, stops, and curtailment events.

\subsection{Thermal and water
management}\label{thermal-and-water-management}

At high current density, the stack generates more heat and requires
adequate water supply and cooling. In a direct DC architecture, the AFE
and the PEM balance of plant must coordinate closely so the stack is
only energised when water flow, temperature, pressure, and conductivity
are within safe limits.

\subsection{Fault handling}\label{fault-handling}

The most critical event is a PEM stack trip while the turbine is
producing power. Without a DC/DC converter, there is less electrical
decoupling between the turbine and the stack. The DC link can
overvoltage quickly unless excess energy is removed or turbine power is
reduced.

A credible direct DC architecture therefore needs:

\begin{itemize}
\tightlist
\item
  fast torque reduction by the AFE;
\item
  pitch or aerodynamic curtailment;
\item
  DC brake chopper or dump load;
\item
  high-speed stack isolation;
\item
  controlled DC-link discharge;
\item
  safe hydrogen shutdown sequence.
\end{itemize}

\section{Implication for the model}\label{implication-for-the-model}

For modelling purposes, direct DC integration can be represented as an
architecture that removes the grid-side inverter and the
electrolyser-side rectifier/DC conversion stage from the
hydrogen-production power path.

The assumption is justified because:

\begin{enumerate}
\def\labelenumi{\arabic{enumi}.}
\tightlist
\item
  the PEM stack is a DC load;
\item
  the turbine already has a controlled DC link;
\item
  industrial concepts explicitly propose direct connection of the
  electrolyser to the turbine DC link;
\item
  academic work supports direct DC bus connection with modular PEM stack
  arrays;
\item
  small-scale experiments have shown direct wind-to-PEM coupling without
  voltage converters;
\item
  the missing DC/DC functions can be replaced by AFE control, stack
  segmentation, and protection systems.
\end{enumerate}

The model should nevertheless treat this as an \textbf{integrated
turbine-electrolyser configuration}, not as a conventional turbine plus
conventional electrolyser package. The feasibility depends on co-design.

\section{Final assessment}\label{final-assessment}

Direct integration of a PEM stack or PEM stack array on the turbine DC
link is a reasonable model assumption for a decentralised hydrogen
architecture.

The main reason to use it is that it avoids unnecessary AC conversion
when the final product is hydrogen rather than electricity. It can
reduce conversion stages, electrical infrastructure, equipment mass, and
conversion losses. It also fits the system logic of a decentralised
hydrogen wind farm, where the infield network collects hydrogen instead
of electricity.

Technically, the assumption is defensible because a PEM stack is a DC
load and a Type IV wind turbine already has a controlled DC link. With
an AFE controlling DC-link voltage and generator torque, and with a
voltage-matched segmented PEM stack array, direct operation without a
separate electrolyser-side DC/DC converter should be possible.

The strongest practical formulation is:

\begin{Shaded}
\begin{Highlighting}[]
\NormalTok{AFE{-}controlled turbine DC link + voltage{-}matched segmented PEM stack array + module switching + emergency dump/brake protection}
\end{Highlighting}
\end{Shaded}

\begin{tcolorbox}[enhanced jigsaw, arc=.35mm, bottomrule=.15mm, bottomtitle=1mm, breakable, colback=white, colbacktitle=quarto-callout-important-color!10!white, colframe=quarto-callout-important-color-frame, coltitle=black, left=2mm, leftrule=.75mm, opacityback=0, opacitybacktitle=0.6, rightrule=.15mm, title=\textcolor{quarto-callout-important-color}{\faExclamation}\hspace{0.5em}{Feasibility statement}, titlerule=0mm, toprule=.15mm, toptitle=1mm]

The model assumption is credible: \textbf{direct DC integration without
a separate electrolyser-side DC/DC converter should be technically
possible}, provided the turbine AFE, PEM stack array, protection system,
and balance of plant are co-designed. It is not yet a widely proven
MW-scale commercial standard, so it should be treated as an advanced
integrated architecture rather than an off-the-shelf configuration.

\end{tcolorbox}




\end{document}
